\section{附录：符号与术语表}\label{sec:appendix}

\begin{table}[htbp]
    \centering
    \caption{本文使用的主要符号}
    \label{tab:notation}
    \begin{tabular}{ll}
        \toprule
        符号 & 含义 \\
        \midrule
        $\zeta_p$ & $p$ 次本原单位根 $e^{2\pi i/p}$ \\
        $\mathbb{Z}[\zeta_p]$ & 分圆整数环 \\
        $h_p$ & 分圆域 $\mathbb{Q}(\zeta_p)$ 的类数 \\
        $E$ & 椭圆曲线 \\
        $L(E, s)$ & 椭圆曲线的 Hasse--Weil $L$ 函数 \\
        $f$ & 权 $2$ 的尖点模形式（新形式） \\
        $\rho_{E,\ell}$ & 椭圆曲线的 $\ell$-adic 伽罗瓦表示 \\
        $G_{\mathbb{Q}}$ & 绝对伽罗瓦群 $\operatorname{Gal}(\overline{\mathbb{Q}}/\mathbb{Q})$ \\
        $R$ & 形变环 \\
        $T$ & Hecke 代数 \\
        $R = T$ & 模性提升定理的代数表述 \\
        $X_0(N)$ & 模曲线（级 $\Gamma_0(N)$） \\
        $P_r$ & 至多 $r$ 个素因子的整数 \\
        \bottomrule
    \end{tabular}
\end{table}

\section*{术语对照}
\begin{itemize}
    \item 无穷下降法（infinite descent）：构造递减解序列导致矛盾的经典反证方法；
    \item 正规素数（regular prime）：不整除分圆域类数的奇素数；
    \item 第一/第二情形（first/second case）：$p \nmid xyz$ 与 $p \mid xyz$ 的两种情形；
    \item 弗雷曲线（Frey curve）：由 FLT 反例构造的半稳定椭圆曲线；
    \item 模性提升（modularity lifting）：由模 $p$ 模性提升到 $\ell$-adic 模性的定理；
    \item 水平约化（level lowering）：将模表示的级降低到其“真实”水平的技巧。
\end{itemize}
